% Required packages in the document preamble:
% \usepackage{booktabs}
% \usepackage{graphicx}

\section{Implementation}
\label{sec:implementation}

The planner is implemented in C++ and operates directly on the generic
environment object produced by the provided parser.  The supplied
\texttt{main} function was left unchanged.  No action names, predicate names,
or domain-specific transition rules are embedded in the planner; all successor
states are derived from the action schemas contained in the input environment.

\subsection{State and Action Representation}

A search state is represented as a set of positive grounded conditions.  We
use the closed-world convention: a condition not contained in the state is
false.  Therefore, a positive precondition is satisfied when its grounded atom
is present, while a negative precondition is satisfied when the corresponding
positive atom is absent.  Applying an action first removes its negative effects
and then inserts its positive effects.  A state satisfies the goal when every
grounded goal condition is present.

Each action schema is grounded by assigning environment symbols to its formal
parameters and substituting those assignments into its preconditions and
effects.  Predicates that never appear in an action effect are static.  During
grounding, an action instance is discarded if one of its static preconditions
is false in the initial state.  Satisfied static preconditions are removed from
the runtime applicability test because their values cannot change.  This
preprocessing is domain independent and cannot remove an action that could
become applicable later.

The state container is unordered, so a deterministic key is constructed for
duplicate detection by sorting the string representation of every grounded
condition.  Search nodes store the state, path cost, parent-node index, and the
action used to reach the node.  Once a goal is found, the plan is reconstructed
by following parent indices back to the initial state.

\subsection{Graph-Search Infrastructure}

Breadth-first search uses a FIFO frontier and marks a state visited when it is
inserted.  This prevents duplicate frontier entries and guarantees a
minimum-action solution for unit-cost actions.  A$^\ast$ uses a priority queue
ordered by $f(s)=g(s)+h(s)$ and stores the lowest discovered $g$-value for each
state.  Stale priority-queue entries are ignored, and a state is reinserted if a
lower-cost path is discovered.  BFS and both A$^\ast$ variants share the same
grounding, applicability, successor-generation, goal-test, state-key, and plan-
reconstruction functions.

The planner records the plan length, grounded-action count, expanded-state
count, generated-successor count, duplicate-successor count, maximum frontier
size, grounding time, search time, and total planning time.  These measurements
are printed before the returned plan.

\subsection{Completeness and Plan Quality}

The reachable state graph is finite because each state is a subset of the
finite set of grounded conditions.  BFS graph search is therefore complete, and
because every action has unit cost, it returns a plan with the minimum number of
actions.  A$^\ast$ with exact empty-delete-list cost $h^+$ is also complete and
optimal for the positive-precondition STRIPS problems used here.  Removing
delete effects produces a relaxation, so its optimal solution cost cannot
exceed the true remaining cost and is therefore admissible.  The normalized
goal-count heuristic is also admissible.  If $M$ is the largest number of goal
literals that any one action can add, then any plan satisfying $k$ currently
unsatisfied goals must contain at least $\lceil k/M \rceil$ actions.  Both
admissible A$^\ast$ variants therefore retain the minimum-action guarantee.

\section{Compilation and Execution}
\label{sec:compilation}

The implementation uses only the C++ standard library and requires a compiler
with C++11 support.  CMake version 3.16.3 or newer may be used, but no external
libraries are required.  From the submitted \texttt{code} directory, the CMake
workflow is:

\begin{verbatim}
mkdir build
cd build
cmake ..
cmake --build . --config Release
\end{verbatim}

On Linux, the executable is normally \texttt{build/planner}.  With a multi-
configuration generator such as Visual Studio, it may instead be
\texttt{build/Release/planner.exe}.  The planner can also be compiled directly
from inside the \texttt{build} directory:

\begin{verbatim}
g++ ../src/planner.cpp -std=c++11 -O2 -o planner
\end{verbatim}

The program takes the basename of an environment file stored in
\texttt{code/envs}.  For example:

\begin{verbatim}
./planner Blocks.txt
\end{verbatim}

The search algorithm is selected with the \texttt{PLANNER\_MODE} environment
variable.  The supported values are \texttt{bfs}, \texttt{astar-goal}, and
\texttt{astar-edl}.  If the variable is not defined, exact EDL A$^\ast$ is used.
Example Linux commands are:

\begin{verbatim}
PLANNER_MODE=bfs        ./planner Blocks.txt
PLANNER_MODE=astar-goal ./planner Blocks.txt
PLANNER_MODE=astar-edl  ./planner Blocks.txt
\end{verbatim}

The corresponding PowerShell syntax is:

\begin{verbatim}
$env:PLANNER_MODE = 'astar-edl'
.\build\Release\planner.exe Blocks.txt
\end{verbatim}

\section{Search Algorithms and Experimental Evaluation}
\label{sec:planner-evaluation}

\subsection{Planner Implementations}

We implemented three forward state-space graph-search planners over the supplied
STRIPS representation.  All actions have unit cost, so plan quality is measured
by the number of actions in the returned plan.

\paragraph{Breadth-first search (BFS).}
BFS serves as the uninformed baseline.  States are marked visited when they are
inserted into the frontier, preventing repeated expansion of the same symbolic
state.  Since every action has unit cost, BFS is complete on the finite state
spaces considered here and returns a minimum-length plan.

\paragraph{A$^\ast$ with goal-count heuristic.}
The first informed planner uses
\begin{equation}
    h_{\mathrm{goal}}(s)
    = \left\lceil
      \frac{\left|\left\{g \in G \mid g \notin s\right\}\right|}{M}
      \right\rceil,
\end{equation}
where $M$ is the maximum number of goal literals appearing in the positive
effects of any one grounded action.  A plan of $n$ actions can add at most $nM$
goal literals, so satisfying $k$ currently unsatisfied goals requires at least
$\lceil k/M \rceil$ actions.  This makes the heuristic admissible while keeping
it inexpensive and domain independent.  In every benchmark domain, each action
adds at most one goal literal, so $M=1$ and the normalized value equals the raw
unsatisfied-goal count.

\paragraph{A$^\ast$ with empty-delete-list heuristic.}
The second informed planner uses the exact empty-delete-list heuristic $h^+$.
For each state evaluated by the outer A$^\ast$ search, a relaxed planning problem
is constructed by removing every negative action effect while retaining action
preconditions.  The heuristic value is the length of the shortest plan in this
relaxed problem.  We solve the relaxed problem with an inner A$^\ast$ search.  An
event-driven $h_{\max}$ propagation supplies an admissible lower bound to the
inner search, but the value returned to the outer search is the exact relaxed
plan length rather than the $h_{\max}$ approximation.  For the positive-
precondition STRIPS domains used in this assignment, $h^+$ is admissible because
removing delete effects cannot make the relaxed problem harder than the original
problem.

\paragraph{Domain-independent preprocessing.}
Before search, action schemas are grounded using the symbols in the environment.
Predicates that never appear in any action effect are identified as static.  A
grounded action is discarded if one of its static preconditions is false in the
initial state, and satisfied static preconditions are removed from its runtime
applicability test.  This pruning is semantics preserving because a static fact
cannot change during planning.  For example, it reduces the large Robot Courier
instance from 6,591 unfiltered symbol assignments to 50 potentially applicable
grounded actions.

\subsection{Experimental Methodology}

The planner was compiled as C++14 with optimization level \texttt{-O2}.  Each
planner/environment combination was executed three times, and the tables report
the median total wall-clock planning time.  Total time includes action grounding,
heuristic initialization, and graph search, but excludes environment-file parsing
and console output.  We additionally recorded the number of expanded states,
generated successors, and maximum frontier size.  Timing measurements were made
under WSL on the same computer for every configuration.  Because sub-millisecond
measurements are sensitive to system noise, expansion counts are the more stable
comparison for the smallest instances.

\subsection{Required Environments}

Table~\ref{tab:required-results} gives results for the supplied environments.
The trivial BlocksEasy instance is included as a correctness check even though
it is not required in the final comparison.  All three planners returned plans
of identical length on every problem.

\begin{table}[t]
    \centering
    \caption{Results on supplied environments. Time is the median of three runs
    and includes both grounding and search.}
    \label{tab:required-results}
    \resizebox{\textwidth}{!}{%
    \begin{tabular}{llrrrrrr}
        \toprule
        Environment & Planner & Grounded & Plan & Expanded & Generated & Max frontier & Time (ms) \\
        \midrule
        BlocksEasy & BFS & 45 & 1 & 5 & 18 & 9 & 0.544 \\
        BlocksEasy & A$^\ast$ goal count & 45 & 1 & 1 & 5 & 5 & 0.406 \\
        BlocksEasy & A$^\ast$ EDL & 45 & 1 & 1 & 5 & 5 & 0.958 \\
        \midrule
        Blocks & BFS & 45 & 3 & 28 & 97 & 15 & 1.152 \\
        Blocks & A$^\ast$ goal count & 45 & 3 & 5 & 26 & 15 & 0.737 \\
        Blocks & A$^\ast$ EDL & 45 & 3 & 3 & 19 & 15 & 8.709 \\
        \midrule
        Blocks and Triangles & BFS & 315 & 6 & 1564 & 15605 & 2197 & 177.794 \\
        Blocks and Triangles & A$^\ast$ goal count & 315 & 6 & 94 & 647 & 237 & 11.818 \\
        Blocks and Triangles & A$^\ast$ EDL & 315 & 6 & 15 & 140 & 95 & 567.683 \\
        \midrule
        Fire Extinguisher & BFS & 117 & 21 & 379 & 2808 & 55 & 14.713 \\
        Fire Extinguisher & A$^\ast$ goal count & 117 & 21 & 336 & 2494 & 55 & 13.539 \\
        Fire Extinguisher & A$^\ast$ EDL & 117 & 21 & 252 & 1869 & 79 & 408.825 \\
        \bottomrule
    \end{tabular}}
\end{table}

The heuristic planners substantially reduce search effort on the nontrivial
Blocks problems.  On Blocks and Triangles, goal-count A$^\ast$ reduces expansions
from 1,564 to 94 (94.0\%), while exact EDL reduces them to 15 (99.0\%).  Despite
this reduction, EDL takes longer than BFS because each outer state evaluation
requires solving an additional relaxed planning problem.  Goal-count A$^\ast$
provides the best runtime on this instance because its heuristic is inexpensive
and still gives strong guidance.

The Fire Extinguisher domain exhibits a smaller heuristic advantage.  Goal-count
A$^\ast$ reduces expansions by 11.3\% relative to BFS, while EDL reduces them by
33.5\%.  Exact EDL is again slower in wall-clock time.  This demonstrates that
the number of expanded states alone does not measure total planning cost: a more
informed heuristic is only beneficial in runtime when the saved expansions
outweigh the cost of computing it.

\subsection{Towers of Hanoi Scaling Study}

Towers of Hanoi was used to test the planners as problem size increases.  The
optimal solution for $n$ disks has $2^n-1$ actions; every planner returned this
known optimum for all four instances.  Table~\ref{tab:hanoi-results} shows the
expected rapid growth in both solution length and planning effort.

\begin{table}[t]
    \centering
    \caption{Towers of Hanoi scaling results.}
    \label{tab:hanoi-results}
    \resizebox{\textwidth}{!}{%
    \begin{tabular}{clrrrrrr}
        \toprule
        Disks & Planner & Grounded & Plan & Expanded & Generated & Max frontier & Time (ms) \\
        \midrule
        3 & BFS & 72 & 7 & 25 & 73 & 8 & 1.689 \\
        3 & A$^\ast$ goal count & 72 & 7 & 14 & 41 & 6 & 1.997 \\
        3 & A$^\ast$ EDL & 72 & 7 & 13 & 38 & 4 & 7.554 \\
        \midrule
        4 & BFS & 126 & 15 & 67 & 200 & 16 & 6.745 \\
        4 & A$^\ast$ goal count & 126 & 15 & 49 & 146 & 10 & 5.053 \\
        4 & A$^\ast$ EDL & 126 & 15 & 45 & 134 & 8 & 88.475 \\
        \midrule
        5 & BFS & 200 & 31 & 237 & 709 & 32 & 23.265 \\
        5 & A$^\ast$ goal count & 200 & 31 & 168 & 503 & 17 & 26.560 \\
        5 & A$^\ast$ EDL & 200 & 31 & 159 & 476 & 16 & 934.404 \\
        \midrule
        6 & BFS & 297 & 63 & 675 & 2024 & 64 & 62.809 \\
        6 & A$^\ast$ goal count & 297 & 63 & 563 & 1688 & 31 & 55.395 \\
        6 & A$^\ast$ EDL & 297 & 63 & 541 & 1622 & 32 & 17703.100 \\
        \bottomrule
    \end{tabular}}
\end{table}

EDL consistently expands the fewest Hanoi states, but its runtime scales poorly.
The delete relaxation does not represent the need to move smaller disks away and
later restore them: once a relaxed action establishes a fact, that fact remains
true.  Consequently, the relaxed problem supplies only a modest reduction in
outer expansions while exact computation of $h^+$ becomes increasingly costly.
For six disks, EDL reduces expansions by 19.9\% relative to BFS but requires
17.7 seconds, compared with 62.8 milliseconds for BFS.  It nevertheless remains
below the assignment's 30-second planning limit.

\subsection{Robot Courier Custom Domain}

The Robot Courier domain tests generality beyond stacking problems.  A mobile
robot navigates a graph of locations, can carry at most one package, and must
deliver each package to its package-specific destination.  Static predicates
encode object types, graph connections, and destinations; dynamic predicates
encode robot/package locations, carrying state, and completed deliveries.  The
large instance increases the problem from four to eight locations and from two
to four packages without changing the action schemas.

\begin{table}[t]
    \centering
    \caption{Results on the custom Robot Courier domain.}
    \label{tab:courier-results}
    \resizebox{\textwidth}{!}{%
    \begin{tabular}{llrrrrrr}
        \toprule
        Instance & Planner & Grounded & Plan & Expanded & Generated & Max frontier & Time (ms) \\
        \midrule
        Courier & BFS & 16 & 10 & 28 & 50 & 6 & 2.060 \\
        Courier & A$^\ast$ goal count & 16 & 10 & 24 & 44 & 5 & 2.809 \\
        Courier & A$^\ast$ EDL & 16 & 10 & 13 & 26 & 6 & 7.737 \\
        \midrule
        Courier Large & BFS & 50 & 30 & 362 & 698 & 22 & 25.882 \\
        Courier Large & A$^\ast$ goal count & 50 & 30 & 342 & 659 & 25 & 27.383 \\
        Courier Large & A$^\ast$ EDL & 50 & 30 & 206 & 401 & 27 & 1268.170 \\
        \bottomrule
    \end{tabular}}
\end{table}

All planners return the same 10-action and 30-action plans for the small and
large instances, respectively.  On Courier Large, EDL reduces expansions by
43.1\% relative to BFS, showing that the relaxed problem provides useful search
guidance.  Its 1.27-second runtime is nevertheless higher than the approximately
26-millisecond uninformed baseline because exact relaxed planning is performed
for many outer states.

\subsection{Discussion}

Across all environments, heuristic search reduces the number of expanded states.
The inexpensive normalized goal-count heuristic provides the best wall-clock
runtime on most instances while retaining admissibility.  Exact EDL is the most
effective method by expansion count, but its nested relaxed searches introduce
a large computational cost.  The results therefore illustrate a central
tradeoff in heuristic planning: heuristic accuracy must be balanced against
evaluation cost.

Static-precondition pruning is important to making the EDL comparison fair.  It
does not alter the relaxed problem or the returned heuristic value, but avoids
considering ground actions that can never become applicable.  Likewise, the
inner $h_{\max}$ computation only orders the exact relaxed search and does not
replace $h^+$.  Thus, the optimized EDL implementation remains faithful to the
empty-delete-list definition while avoiding unnecessary work.
